\minitopic{比较研究包含三种不同事业}历史解释力求说明一个概念或论证在原语言、文本位置和实践中的意义；结构比较选择共同问题轴，考察不同材料怎样处理证据、错误、能力或行动；规范建构则借多种资源提出今天可辩护的新理论。三者可以连续进行，却必须在写作中标明转换。若把自己的重构写成古代作者原意，会伪造历史主张；若拒绝任何重构，又会使历史材料只能被陈列而不能参与当前论证。

\minitopic{比较单位不能只是一对词}“知”与knowledge、“信”与belief可能在部分语境重叠，却各自嵌在对象、推论、角色和实践中。可靠比较应描述一个概念簇：它典型适用于什么，和哪些词对立，允许推出什么，失败案例是什么，谁有资格使用，产生何种行动后果。比较单位也可以是完整论证或问题结构，例如“如何在有限视角中修正判断”，而不是假定两边先有同名概念。单位选得越清楚，结论的强度越容易控制。

\minitopic{内部解释要求文本位置和反方材料}摘取一句“知之为知之”即可支持多种现代美德；只有放回《论语》相应对话和学习、言说规范中，才能判断它主要约束诚实自评、求学态度还是知识条件\parencite[2.17]{analects2003}。同样，庄周梦蝶不能因包含梦就直接成为笛卡尔怀疑论。内部解释要追踪篇章功能、术语变化、角色话语与作者立场，还要寻找不符合现代对应的材料。恰恰是这些阻力，能够防止比较退化为名句装饰。

\minitopic{四项检验控制相似主张}语义检验问关键术语范围是否相近；论证检验问双方是否依相似前提和推理，而非只有结论相似；功能检验问概念在法庭证明、修身、科学解释或政治劝诫中做什么；历史检验问相似是否来自传播、共同文本或独立平行。通过一项不代表通过四项。两个概念可以语义不同却解决相似功能，也可以字面相似却服务完全不同实践。结论应写成“在哪一轴上、到何种程度相似”。

\begin{argumentbox}
每提出一次跨传统对应，附一张限度表：最强相似证据；最强不相似证据；不可确定之处；是否主张历史传播；该对应对双方各改变什么。若对应只让熟悉理论显得更古老，却没有反向修正任何问题，其认识收益通常很低。
\end{argumentbox}

\minitopic{厚概念不能被抽成薄条件}某些“知”同时包含能力、修养、情境判断或行动表现，而命题知识分析可能只问主体是否对真命题具有适当信念。把厚概念还原成薄状态，会丢失规范与实践；反过来，用厚概念批评薄概念“不够完整”，也可能误解后者本来只承担局部理论角色。比较应先承认厚薄差异，再问薄分析是否需要解释行动关联、厚概念是否能区分真值与德性。差异由此成为双向问题，而非胜负证据。

\minitopic{历史同源需要独立证据}年代更早不证明影响，概念相似也不证明传播。历史同源主张至少需要可追踪文本、翻译、接触网络、术语借用或同时代接受记录。没有这些证据时，可以主张独立平行或结构相似，但不能使用“早已提出”“影响了”等强词。谨慎不是削弱比较，而是让哲学结构与历史因果各自接受合适证据。

\minitopic{翻译应保存推论而不只保存字面}一个译词看似自然，却可能让读者自动带入原概念没有的反义词、必要条件与理论争论。判断译法时可做推论测试：在原文中，拥有该状态能推出什么，缺少它是否受责备，它与学习、行动或角色怎样连接；再检查译入语中的概念是否允许同样推论。若不能，应保留原词、使用复合译法或加限定说明。译文的目标不是让陌生材料立刻熟悉，而是让读者既能进入论证，又看见不可无损转换之处。

\minitopic{案例跨语境迁移需要重建背景常识}银行、谷仓或法庭案例依赖特定制度、风险和语言习惯，直译到另一环境后，读者可能因不熟悉支票制度、财产权或证言角色而作不同判断。差异未必显示知识概念不同，也可能显示案例中的默认事实改变。比较实验应先共同确认制度背景，再制作本地等功能案例，并记录哪些变量无法对齐。只有判断差异在多种表达和背景控制后仍稳定，才有理由把它提升为概念或规范差异。

\minitopic{比较关系中的权力会决定谁负责解释}当一种传统被当作默认理论语言，另一种只被要求提供“文化例子”时，比较已经不对称。前者的概念无需说明历史位置，后者却被迫证明能翻译成前者。方法上的修正不是取消共同标准，而是让双方都交代语境、反例与盲点，并允许非默认材料改变问题分类。选取文本、译本和对话对象时还应说明谁的解释被代表，防止把内部多样传统压成一个整齐声音。

\minitopic{谱系批评与规范反驳承担不同任务}说明一个概念如何在殖民、宗教、科学或官僚制度中形成，可以揭示其排除与权力作用，却不单独证明该概念为假；给出逻辑反例也不能说明概念为何在历史中取得权威。比较认识论应把两种批评分工：谱系追踪问题和分类怎样被制造，规范论证检验它们是否值得保留。二者结合时，历史材料可提出被现有理论忽略的代价，规范分析则防止“出身有问题”直接变成否定结论。

\minitopic{负面比较也能产生实质知识}研究可能发现两个概念没有共同对象、两套论证依赖不兼容实践，或相似仅是译词制造。这样的结果不是项目失败。它能阻止错误历史主张，澄清各理论的边界，并指出真正共同问题应上移或下移到何处。负面结论也需证据：展示曾被主张的对应、系统列出不相似材料，并说明为何更弱比较仍不成立。与轻率统一相比，有论证的不可对应更能为后续研究节省道路。

\begin{argumentbox}
比较写作的最小证据包包括：原文与版本；至少一种竞争解释；关键术语的推论表；现代对话理论的最强表述；共同问题轴；相似与不相似的对称证据；历史传播证据状态；以及研究者自己的建构从哪一步开始。
\end{argumentbox}

\minitopic{阶段性结论}比较认识论的首要纪律是控制主张强度。它不从词形相似出发，而从明确问题、内部解释和关系结构出发；不把现代框架设为唯一裁判，也不把传统差异神秘化。经过语义、论证、功能与历史检验后，研究者才能决定应当翻译、保留差异，还是公开进行规范重构。
